\chapter{Самодиффузия в углеводородах}\label{ch:ch2}

В некоторых случаях молекулярной динамике свойственно неточно воспроизводить экспериментальные данные, например, зависимость плотности от температуры и давления от температуры. Однако, при этом зависимость плотности от давления может воспроизводиться точно. Это возможно по ряду причин: в случае некорректной оценки дальнодействующих поправок, особенностей параметризации потенциала межатомного взаимодействия и других тонкостей молекулярного моделирования.

Эта глава посвящена влиянию некоторых из факторов на результат молекулярного моделирования. Изучены сравнительные особенности различных моделей при моделировании углеводородов. Были испробованы модели TraPPE-UA~\cite{trappeua}, OPLS-AA~\cite{jorgensen_development_1996} и COMPASS~\cite{sun_compass_1998}. Они различаются подробностью описания молекулы: TraPPE-UA \--- модель объединенного атома, OPLS-AA \--- стандартная полноатомная модель, COMPASS \--- подробная модель, в которой имеется взаимодействие между связями, углами и двугранными углами.

Рассмотрена диффузия углеводородов моделированием методом молекулярной динамики. Использована формула Эйнштейна-Смолуховского. Рассмотрено 2 вещества: 2,2,4-триметилпентан$(C_8H_{18})$ и 1,1-дифенилэтан$(C_{14}H_{14})$. 

В первом подразделе даны подробные описания этих моделей, показано, как в них учтены различные виды взаимодействий. Далее описан способ вычисления коэффициента диффузии и полученные результаты. В последнем подразделе приведены краткие выводы из главы. Результаты, приведенные в данной главе, опубликованы в~\cite{kondratyuk_transport_2020}.

\section{Используемые потенциалы межатомного взаимодействия}

В работе используется три модели различной сложности, полноатомные OPLS-AA(Optimized Potentials for Liquid Simulations \-- All Atoms) и COMPASS(Condensed-phase Optimized Molecular Potentials for Atomistic Simulation Studies) и модель объединенного атома TraPPE-UA(Transferable Potential for Phase Equilibrium \-- United Atom). 

Потенциальная энергия рассчитывается по формуле:

\begin{align}
	U = U_{bond} + U_{angle} + U_{dihedral} + U_{improper} + U_{lj} + U_{coul} ,
\end{align}

где первые четыре слагаемых \--- внутримолекулярная энергия: энергия связей~$U_{bond}$, углов~$U_{angle}$, двугранных углов~$U_{dihedral}$ и квадрупольного взаимодействия~$U_{improper}$. Два последних слагаемых относятся к межмолекулярному взаимодействию, это $U_{lj}$ и $U_{coul}$ соответственно. 

Связи и углы в системе описывались с помощью гармонических осцилляторов в случае OPLS-AA и TraPPE-UA:

\begin{align}
	U_{bond} = k (r - r_0)^2 ,
\end{align}

где $k$ \--- константа жесткости, $r$ \--- длина связи, $r_0$ \--- равновесная длина связи.

В случае COMPASS для тех же случаев использовался негармонический осциллятор:

\begin{align}
	U_{bond} = k_1 (r - r_0)^2 + k_2 (r - r_0)^2 + k_3 (r - r_0)^2,
\end{align}

где $k_1, k_2, k_3$ \--- константы.

Также для COMPASS имеется энергетическая взаимосвязь между связями и углами:

\begin{subequations}
    \begin{gather}
    		U_{angle} = U_a + U_{bb} + U_{ba} \\
		U_{a} = k_1 (\theta - \theta_0)^2 + k_2 (\theta - \theta_0)^2 + k_3 (\theta - \theta_0)^2 \\
		U_{bb} = M (r_{ij} - r_1 )(r_{jk} - r_2 ) \\
		U_{ba} = N_1 (r_{ij} - r_1 )(\theta - \theta_0 ) + N_2 (r_{jk} - r_2)(\theta - \theta_0)
    \end{gather}
\end{subequations}

где $\theta_0$ \--- равновесный угол, $r_1$ и $r_2$ \--- равновесные длины связей, образующих угол.

Для двугранных углов во всех моделях использовались суммы Фурье. Они выглядят следующим образом:

\begin{align}
		U = k_1 [1 - cos(\phi - \phi_1 )] + k_2 [1 - cos(2\phi - \phi_2 )] + k_3 [1 - cos(3\phi - \phi_3 )],
\end{align}


где $\phi_1$, $\phi_2$, $\phi_3$ \--- равновесные углы, $k_1$, $k_2$, $k_3$ \--- константы. Как и с углами, для COMPASS учитывается взаимодействие между связями, углами и двугранными углами:

\begin{subequations}
    \begin{gather}
    		U_{dihedral} = U_d + U_{mbt} + U_{ebt} + U_{at} + U_{aat} + U_{bb13} \\
		U_d = k_1 [1 - cos(\phi - \phi_1 )] + k_2 [1 - cos(2\phi - \phi_2 )] + k_3 [1 - cos(3\phi - \phi_3 )] \\
		U_{mbt} =(r_{jk} - r_2 )[A_1 cos(\phi) + A_2 cos(2\phi) + A_3 cos(3\phi)] \\
		U_{ebt} =(r_{ij} - r_1)[B_1 cos(\phi) + B_2 cos(2\phi) + B_3 cos(3\phi)]+(r_{kl} - r_3 )[C_1 cos(\phi) + C_2 cos(2\phi) + C_3 cos(3\phi)] \\
		U_{at} =(\theta_{ijk} - \theta_1)[D_1 cos(\phi) + D_2 cos(2\phi) + D_3 cos(3\phi)]+(\theta_{jkl} - \theta_2 )[E_1 cos(\phi) + E_2 cos(2\phi) + E_3 cos(3\phi)] \\
		U_{aat} = M (\theta_{ijk} - \theta_1)(\theta_{jkl} - \theta_2) cos(\phi) \\
		U_{bb13} =N (r_{ij} - r_1)(r_{kl} - r_3)
    \end{gather}
\end{subequations}
 
где $r_1$, $r_2$, $r_3$ \--- равновесные длины связей.

Межмолекулярные взаимодействия описываются потенциалом Леннард-Джонса. В случае TraPPE-UA и OPLS-AA использован вариант со степенями 12-6, для COMPASS \--- 9-6, и такой вариант лучше описывает отталкивание на малых дистанциях. Радиус обрезания для полноатомных моделей равен 12~\AA, для TraPPE-UA \--- 14~\AA. Кулоновское взаимодействие описывается методом particle-particle particle-mesh, в котором заряды проецируются на решетку.

Использована иерархическая схема rRESPA, в которой межмолекулярные взаимодействия рассчитаны с шагом 2~фс, связи и углы \--- с шагом 0.5~фс, двугранные углы \--- 1~фс. Расчет проводился с периодическими граничными условиями.

Начальная конфигурация состояла из 1000~молекул $C_8H_{18}$ и 729~молекул $С_{14}H_{14}$. Молекулы $C_8H_{18}$ расположены случайным образом и находятся при плотности 0.7~г/см$^3$. Для $С_{14}H_{14}$ начальным состоянием является идеальный газ с плотностью примерно 1 г/см$^3$. Первой стадией моделирования является минимизация энергии. Далее идет вывод на равновесие сначала в NPT-ансамбле в течение 1~нс, затем в NVT также в течение 1~нс. Используется термостат Нозе-Гувера с временем релаксации 100~фс и в случае NPT-ансамбля \--- баростат Шиноды. 

По равновесной части NPT-расчета проводится усреднение плотности. Для аппроксимации используется уравнение~\cite{dymond_tait_1988}:

\begin{align}
		(\rho - \rho_0)/\rho = A*log_{10}[(B+P)/(B+P_0)],
\end{align}

где $\rho$ и $P$ \--- плотность и давление, $\rho_0$ и $P_0$ \--- плотность и давление при 0.1~МПа, $A$ и $B$ \--- константы, определяемые через аппроксимацию.

\begin{figure}[ht]
    \centerfloat{
        \includegraphics[scale=0.75]{densities.png}
    \caption{Зависимости плотности от давления для $C_8H_{18}$, полученные с помощью моделей TraPPE-UA(треугольники), OPLS-AA(квадраты) и COMPASS(круги), а также в эксперименте(ромбы).}
    }\label{fig:densities_C8}
\end{figure}

\begin{figure}[ht]
    \centerfloat{
        \includegraphics[scale=0.75]{densities.png}
    \caption{Зависимости плотности от давления для $C_{14}H_{14}$, полученные с помощью модели COMPASS(круги) и в эксперименте(ромбы).}
    }\label{fig:densities_C14}
\end{figure}

Результаты для всех моделей примерно совпадают и лежат в пределах 5 процентов от экспериментальных~\cite{schiessler_properties_1955, dymond_transport_1980}, несмотря на различие в подробности описания молекул. В случае с COMPASS такой результат ожидаем в силу сложности модели. TraPPE-UA была параметризована под условия, подобные исследуемым, а также лучше учитывает изменения электронной плотности в зависимости от удаленности атома углерода от конца цепи. OPLS-AA дает завышенные результаты относительно обеих моделей и эксперимента. Это происходит из-за единообразного описания всех групп $CH_2$ независимо от их удаленности от конца цепи. Расхождение с экспериментом возрастает с ростом давления. Таким образом, в зависимости от модели можно предсказать особенности описания ей вещества в выбранных условиях. А с помощью этих данных можно ввести поправки, которые помогут привести данные в соответствие с экспериментальными.

\section{Коэффициент диффузии}

Для вычисления коэффициента диффузии использовалась формула Эйнштейна-Смолуховского:

\begin{align}
		\underset{t \to \infty}{\lim} \langle \Delta r(t)^2 \rangle = 6Dt + const,
\end{align}

где $\langle \Delta r \rangle$ \--- среднее смещение атомов от начального положения в момент времени $t$, $D$ \--- коэффициент диффузии. Коэффициент определяется по линейной части графика средневадратичного смещения от времени. 

\begin{figure}[ht]
    \centerfloat{
        \includegraphics[scale=0.75]{diffusion.png}
    \caption{Зависимости коэффициентов диффузии от давления для $C_8H_18$, полученные с помощью моделей TraPPE-UA(треугольники), OPLS-AA(квадраты) и COMPASS(круги). Пунктирными линиями показана аппроксимация по формуле~\ref{eq:diffusion}.}
    }\label{fig:diffusion_C8}
\end{figure}
\begin{figure}[ht]
    \centerfloat{
        \includegraphics[scale=0.75]{diffusion_C14.png}
    \caption{Зависимости коэффициентов диффузии от давления для $C_14H_14$, полученные с помощью модели COMPASS(круги). Пунктирными линиями показана аппроксимация по формуле~\ref{eq:diffusion}.}
    }\label{fig:diffusion_C14}
\end{figure}
Для модели TraPPE-UA коэффициент диффузии оказался переоценен по сравнению с COMPASS. Вероятно, это связано с большим количеством свободного объема за счет объединения атомов водорода с атомами углерода. Модель OPLS-AA занижает значения коэффициента диффузии, поскольку плотность в этой модели оказалась завышенной по сравнению с другими и экспериментом. Таким образом, по поведению изотермы можно предсказать зависимость коэффициента диффузии от давления. Это означает, что по полученной из экспериментальных данных плотности и дальнейшему молекулярному моделированию системы с такой плотностью можно предсказать коэффициент самодиффузии, который не был определен в эксперименте. 

Зависимость коэффициента диффузии от давления аппрокисимировалась следующим уравнением:

\begin{align}
		\label{eq:diffusion}
		ln(D(P)/D_0) = K P + L P^2,
\end{align}

где $D_0$ \--- коэффициент диффузии при 0.1~МПа, $K$ и $L$ \--- константы. Видно, что данная зависимость хорошо описывает полученные данные.
 

\section{Выводы}

Проведены вычисления диффузии методом молекулярной динамики для 2,2,4-триметилпентана и 1,1-дифенилэтана с помощью формулы Эйнштейна-Смолуховского. Были использованы 3 модели: TraPPE-UA, OPLS-AA и COMPASS. Расчёт проводился для диапазона давлений от 0.1 до 400~МПа. Температура была постоянной и равной 298~К. 

Было проанализировано влияние выбора модели для описания углеводородов на их плотность. Показано, что для различных моделей можно предсказать, как они повлияют на конечный результат. Таким образом возможна коррекция полученных для более легкой модели результатов так, чтобы они совпадали с экспериментом или с более сложной моделью. 

Для коэффициента диффузии референтными стали значения, полученные с помощью модели COMPASS. Значения, полученные с помощью модели TraPPE-UA, оказались завышенными в силу более низких значений плотности. Это означает, что зависимость коэффициента диффузии от плотности воспроизводится, в отличие от зависимости коэффициента диффузии от давления или плотности от давления. Это позволяет получить коэффициент самодиффузии при нужном давлении, сделав поправку к плотности, и лишает необходимости ставить дополнительный эксперимент или усложнять модель. Аналогичный эффект свойственен модели OPLS-AA. Поскольку эта модель завышает плотность, для неё полученные значения коэффициента самодиффузии будут заниженными. Соответственно, поправку к плотности потребуется сделать со знаком противоположным этой поправке в случае модели TraPPE-UA.   

Поскольку нет референсных результатов для диффузии, как таковые использовались значения диффузии для модели COMPASS, так как её значения плотности были ближе всего к экспериментальным. Значения, полученные с помощью модели TraPPE-UA, оказались завышенными, что объясняется недооценкой плотности в этой модели. Для OPLS-AA наоборот, вычисленная диффузия получилась ниже референтной, что является следствием большей плотности, чем в модели COMPASS.    